\chapter{The Problem}
\label{ch:the-problem}

\textit{Every day, millions of developers sit down with an AI assistant and make
dozens of consequential decisions---choosing databases, designing APIs,
selecting architectures.  The code that results from these conversations
survives.  The reasoning that produced it does not.}

\section{AI-Assisted Development}
\label{sec:ai-dev}

Software development in 2025 looks fundamentally different from even two years
ago.  Over 90\% of professional developers report using AI coding assistants at
least weekly, and tools like GitHub Copilot, Cursor, and Claude Code have
become as ubiquitous as version control itself.  A typical coding session now
involves dozens of conversational exchanges: the developer describes intent,
the assistant proposes options, and together they iterate toward a solution.

Each of these exchanges carries information that traditional software
engineering artifacts were never designed to capture.  A single afternoon
session might involve:

\begin{itemize}
    \item Evaluating three caching strategies before choosing Redis with a
          TTL-based invalidation scheme.
    \item Deciding to use PostgreSQL over MongoDB after discussing ACID
          requirements.
    \item Rejecting a microservices split in favor of a modular monolith after
          weighing operational complexity.
    \item Choosing FastAPI over Flask because of native async support and
          automatic OpenAPI generation.
\end{itemize}

These are not trivial preferences---they are \emph{architectural decisions}
with months-long consequences.  And they happen entirely inside a conversation
window.

\begin{keyinsight}
AI-assisted development has made decision-making more explicit than ever.
Developers literally \emph{write out} their reasoning when prompting an
assistant.  The problem is not that rationale is missing---it is that we throw
it away.
\end{keyinsight}

\section{The Rationale Gap}
\label{sec:rationale-gap}

When the conversation window closes, the rationale evaporates.  What remains
is code: functions, configuration files, dependency manifests.  The code tells
you \emph{what} was built, but not \emph{why}.

Three months later, a new team member opens a pull request and asks:

\begin{quote}
``Why did we choose PostgreSQL over MongoDB for the user service?  The access
patterns look document-oriented.''
\end{quote}

Nobody remembers.  The original developer has moved to another project.  The
conversation where the trade-offs were weighed is buried in a chat log that
nobody indexes, searches, or even knows exists.  The new developer makes a
reasonable but uninformed decision to migrate---and re-introduces the very
problem the original choice was designed to prevent.

This is the \textbf{rationale gap}: the systematic loss of decision context
that occurs when the medium of decision-making (conversation) is disconnected
from the medium of record (code and documentation).

\begin{warning}
The rationale gap is not merely inconvenient.  It leads to repeated
architectural mistakes, contradictory decisions across teams, and wasted
engineering effort re-deriving conclusions that were already reached.
\end{warning}

\section{Existing Approaches}
\label{sec:existing-approaches}

The software engineering community has tried several approaches to close the
rationale gap.  None have succeeded at scale.

\begin{table}[ht]
\centering
\caption{Existing approaches to capturing decision rationale.}
\label{tab:existing-approaches}
\begin{tabularx}{\textwidth}{l c c c c X}
\toprule
\textbf{Approach} & \textbf{Auto?} & \textbf{Rationale?} & \textbf{Search?} & \textbf{Agent?} & \textbf{Problem} \\
\midrule
ADRs           & No  & Yes       & Partially & No & $<$5\% adoption; manual effort \\
Commit msgs    & No  & Rarely    & Yes       & No & Captures \emph{what}, not \emph{why} \\
Documentation  & No  & Sometimes & Yes       & No & Perpetually stale \\
Slack / Chat   & No  & Yes       & Barely    & No & Ephemeral, unsearchable \\
\bottomrule
\end{tabularx}
\end{table}

\textbf{Architecture Decision Records (ADRs)} are the gold standard in theory.
Each ADR documents a single decision with its context, options considered, and
rationale.  In practice, fewer than 5\% of engineering organizations maintain
ADRs consistently.  They require manual effort at precisely the moment when
developers are least motivated to write documentation---right after making the
decision and wanting to move on to implementation.

\textbf{Commit messages} are automated in the sense that every commit requires
one, but they overwhelmingly describe \emph{what} changed rather than
\emph{why}.

\textbf{Documentation} captures rationale when written well, but documentation
is notoriously difficult to keep synchronized with the codebase.  The half-life
of accurate documentation in a fast-moving project is measured in weeks.

\textbf{Slack and chat logs} often contain the richest rationale, but they are
buried in ephemeral message streams, unsearchable by future developers, and
completely inaccessible to AI agents.

\begin{keyinsight}
Every existing approach either requires manual effort (which kills adoption) or
captures the wrong information (what changed, not why).  The challenge is
finding a source that is both \emph{automatic} and \emph{rationale-rich}.
\end{keyinsight}

\section{The Insight}
\label{sec:the-insight}

The insight behind \sysname{} is deceptively simple:

\begin{quote}
\emph{The AI conversation \textbf{is} the decision record.}
\end{quote}

When a developer discusses options with an AI assistant, the conversation
naturally contains every element of a well-written ADR:

\begin{enumerate}
    \item \textbf{Trigger} --- what prompted the decision.
    \item \textbf{Context} --- background constraints.
    \item \textbf{Options considered} --- alternatives discussed.
    \item \textbf{Decision} --- what was chosen.
    \item \textbf{Rationale} --- why it was chosen.
\end{enumerate}

This information does not need to be \emph{created}---it already exists.  It
just needs to be \emph{extracted} and \emph{structured}.

That is the core thesis of \sysname{}: passively capture AI coding
conversations, extract decision traces using a large language model, resolve
the entities mentioned into a canonical knowledge graph, and make the result
searchable by both humans and AI agents.  No manual effort.  No behavior
change.  The developer keeps working exactly as they already do.

\section{Worked Example}
\label{sec:worked-example}

To make this concrete, consider the following exchange between a developer and
an AI coding assistant (Listing~\ref{lst:conv}).

\begingroup
\lstset{style=bash, numbers=none, basicstyle=\small\ttfamily}
\begin{lstlisting}[caption={Developer--AI conversation about database selection.},label={lst:conv},escapechar=\%]
%\textbf{Developer:}% I need to select a database for a new
  transaction processing service on our fintech platform.
  Requirements: strong consistency, ACID transactions,
  and complex queries on financial records.

%\textbf{Assistant:}% For a transaction-heavy financial service,
  the main contenders would be PostgreSQL, MongoDB, and
  CockroachDB. Let me compare them:
  - PostgreSQL: Mature ACID support, rich SQL, excellent
    for complex joins on financial data.
  - MongoDB: Flexible schema, but multi-document
    transactions add overhead and complexity.
  - CockroachDB: Distributed ACID, but adds operational
    complexity you may not need yet.
  Given your emphasis on ACID compliance and complex
  queries, PostgreSQL is the strongest fit.

%\textbf{Developer:}% Good point about MongoDB transactions.
  We also need PostGIS later for location-based fraud
  detection. Going with PostgreSQL.
\end{lstlisting}
\endgroup

\noindent
From this 20-line exchange, \sysname{} extracts a structured \textbf{decision
trace} (Table~\ref{tab:extracted-trace}).

\begin{table}[ht]
\centering
\caption{Decision trace extracted by \sysname{} from Listing~\ref{lst:conv}.}
\label{tab:extracted-trace}
\begin{tabularx}{\textwidth}{l X}
\toprule
\textbf{Field} & \textbf{Extracted Value} \\
\midrule
Trigger    & Need to select database for new service \\
Context    & Building a transaction-heavy financial API \\
Options    & PostgreSQL, MongoDB, CockroachDB \\
Decision   & Use PostgreSQL \\
Rationale  & ACID compliance critical for financial transactions;
             PostGIS needed for future fraud detection \\
Confidence & 0.92 \\
\bottomrule
\end{tabularx}
\end{table}

Each entity mentioned---PostgreSQL, MongoDB, CockroachDB, PostGIS---is resolved
against \sysname{}'s knowledge graph using the seven-stage entity resolution
pipeline (Chapter~\ref{ch:entity-resolution-pipeline}).  The decision node is
linked to these entities via \ic{INVOLVES} relationships, making it
discoverable through graph traversal.

\begin{keyinsight}
A single short conversation produced a complete decision record---trigger,
context, alternatives, choice, and rationale---without the developer writing
any documentation at all.  This is the power of treating conversations as
first-class engineering artifacts.
\end{keyinsight}
